\documentclass[10pt,reqno]{amsart}

\usepackage[T1]{fontenc}
\usepackage[utf8]{inputenc}
\usepackage{newtxtext,newtxmath}
\usepackage{microtype}
\usepackage{mathtools}
\usepackage{booktabs}
\usepackage{enumitem}
\usepackage{xcolor}
\usepackage{hyperref}
\usepackage[nameinlink,noabbrev]{cleveref}
\usepackage[a4paper,margin=28mm,headheight=15pt]{geometry}

\definecolor{LinkBlue}{HTML}{174A7E}
\hypersetup{
  colorlinks=true,
  linkcolor=LinkBlue,
  citecolor=LinkBlue,
  urlcolor=LinkBlue,
  pdftitle={A Positive-Characteristic Counterexample to the Literal Q-Cartier Base-Change Equivalence},
  pdfauthor={Argus Open Research Project},
  pdfsubject={Weak Q-Gorensteinness and pluricanonical base change; Result 20000320},
  pdfkeywords={Q-Gorenstein morphism, reflexive power, base change, surface family, positive characteristic}
}

\setlength{\parskip}{2pt}
\setlist[enumerate]{leftmargin=7mm,itemsep=2pt,topsep=3pt}
\emergencystretch=2em

\newtheorem{theorem}{Theorem}[section]
\newtheorem{proposition}[theorem]{Proposition}
\newtheorem{lemma}[theorem]{Lemma}
\newtheorem{corollary}[theorem]{Corollary}
\theoremstyle{definition}
\newtheorem{definition}[theorem]{Definition}
\newtheorem{remark}[theorem]{Remark}

\title[A counterexample to a Q-Gorenstein base-change equivalence]
{A Positive-Characteristic Counterexample to the Literal Q-Cartier/Base-Change Equivalence}

\author{Argus Open Research Project}
\address{Argus Open Research Project}
\date{August 30, 2026}

\subjclass[2020]{Primary 14D06; Secondary 14J10, 14B12}
\keywords{$\mathbb Q$-Gorenstein morphism, reflexive pluricanonical sheaf, base change, surface family}

\begin{document}

\begin{abstract}
AIM Problem~1.1 asks whether the condition that the relative dualizing sheaf of a family of surfaces is $\mathbb Q$-Cartier is equivalent to Koll\'ar's requirement that all reflexive pluricanonical powers commute with arbitrary base change. We verify that the unrestricted equivalence is false. Lee and Nakayama construct a flat projective surface family whose first infinitesimal thickening over $B=\operatorname{Spec}k[t]/(t^2)$, in positive characteristic, is naively or weakly $\mathbb Q$-Gorenstein, while its second reflexive pluricanonical power fails relative $S_2$. Their base-change criterion then implies that the map to the reduced closed fiber is not an isomorphism. This gives a direct counterexample to the weak-to-Koll\'ar implication. The example has a nonreduced Artin base and Fano surface fibers; it does not contradict Hacking's characteristic-zero DVR theorem or settle a silently narrowed canonically polarized formulation.
\end{abstract}

\maketitle

\section{The literal equivalence}

Let $f:X\to B$ be a family of surfaces with relative dualizing sheaf $\omega_{X/B}$. For $m>0$, write
\[
  \omega_{X/B}^{[m]}=\bigl(\omega_{X/B}^{\otimes m}\bigr)^{\vee\vee}.
\]
For a base change $\phi:B'\to B$, with induced map $\phi_X:X\times_BB'\to X$, the relevant comparison is
\begin{equation}\label{eq:basechange}
  \bigl(\phi_X^*\omega_{X/B}\bigr)^{[m]}
  \longrightarrow
  \phi_X^*\bigl(\omega_{X/B}^{[m]}\bigr).
\end{equation}
AIM Problem~1.1 asks whether the following two conditions are equivalent \cite{VanOpstall}:
\begin{enumerate}[label=\textup{(\roman*)}]
  \item $\omega_{X/B}$ is $\mathbb Q$-Cartier, equivalently some positive reflexive power is invertible;
  \item for every $m>0$ and every base change $B'\to B$, the comparison \eqref{eq:basechange} is an isomorphism.
\end{enumerate}
The source calls (ii) Koll\'ar's stronger condition. It places no reduced-base or characteristic-zero assumption in the question itself.

Lee and Nakayama use the following terminology \cite[Definition~7.1]{LeeNakayama}.

\begin{definition}
For an $S_2$ morphism with $\mathbb Q$-Gorenstein fibers:
\begin{enumerate}[label=\textup{(\alph*)}]
  \item the morphism is \emph{naively $\mathbb Q$-Gorenstein} if $\omega^{[r]}$ is invertible for some $r>0$;
  \item it is \emph{$\mathbb Q$-Gorenstein} if every $\omega^{[m]}$ satisfies relative $S_2$.
\end{enumerate}
\end{definition}
Their introduction identifies (a) with the weak or Viehweg condition and (b) with the Koll\'ar condition. Along a $\mathbb Q$-Gorenstein fiber, relative $S_2$ for $\omega^{[m]}$ is equivalent to compatibility with the fiber base-change map \cite[Remark~7.6]{LeeNakayama}.

\section{The projective surface family}

Lee and Nakayama begin with the contraction of the $(-4)$-curve on the Hirzebruch surface $F_4$, obtaining the weighted projective plane
\[
  X=\mathbb P(1,1,4).
\]
Their Example~7.4 constructs, over an algebraically closed field $k$, a flat projective morphism
\[
  f:Y\longrightarrow T=\mathbb A^1_k
\]
with the following properties:
\begin{enumerate}
  \item $Y$ is normal and projective;
  \item the special fiber is $Y_0\simeq\mathbb P(1,1,4)$, a $\mathbb Q$-Gorenstein surface of index $2$;
  \item for $t\neq0$, one has $Y_t\simeq\mathbb P^1\times\mathbb P^1$;
  \item every fiber is $\mathbb Q$-Gorenstein.
\end{enumerate}
The original family over $\mathbb A^1$ is not naively $\mathbb Q$-Gorenstein: the canonical self-intersection jumps from $K_{Y_0}^2=9$ to $K_{Y_t}^2=8$.

Localize $T$ at the origin. Let $R$ be the resulting discrete valuation ring, with maximal ideal $\mathfrak m$, and define
\[
  R_n=R/\mathfrak m^{n+1},
  \qquad T_n=\operatorname{Spec}R_n,
  \qquad Y_n=Y\times_TT_n.
\]
In positive residue characteristic, the infinitesimal thickenings behave differently from the whole DVR family.

\section{The Artin-thickening counterexample}

\begin{theorem}\label{thm:counterexample}
Let $k$ be an algebraically closed field of characteristic $p>0$. There exists a flat projective family of surfaces
\[
  g:Y_1\longrightarrow B=\operatorname{Spec}k[t]/(t^2)
\]
such that $g$ is naively $\mathbb Q$-Gorenstein, but $\omega_{Y_1/B}^{[2]}$ fails relative $S_2$. In particular, the closed-point base-change map
\[
  \omega_{Y_1/B}^{[2]}\otimes_B k
  \longrightarrow
  \omega_{\mathbb P(1,1,4)/k}^{[2]}
\]
is not an isomorphism.
\end{theorem}

\begin{proof}
Apply the preceding construction in characteristic $p>0$. Since the special fiber is $\mathbb Q$-Gorenstein, Lee--Nakayama Remark~7.28 implies that every infinitesimal thickening $Y_n\to T_n$ is naively $\mathbb Q$-Gorenstein. For $n=1$, the base is $B=\operatorname{Spec}k[t]/(t^2)$, so some positive reflexive power of $\omega_{Y_1/B}$ is invertible.

For this same first thickening, Lemma~7.5 of \cite{LeeNakayama} states that $\omega_{Y_1/B}^{[2]}$ is not relatively $S_2$ and that $Y_1\to B$ is not $\mathbb Q$-Gorenstein in the stronger sense. Remark~7.6 identifies this relative-$S_2$ failure with failure of the displayed base-change map to the closed fiber. The target is invertible because the special fiber has index $2$.
\end{proof}

\begin{corollary}\label{cor:aim}
The unrestricted equivalence in AIM Problem~1.1 is false: weak or naive $\mathbb Q$-Gorensteinness does not imply compatibility of all reflexive pluricanonical powers with arbitrary base change.
\end{corollary}

\begin{proof}
The family in \Cref{thm:counterexample} satisfies condition (i) of the AIM question and fails condition (ii) for $m=2$ and the base change from $B$ to its reduced closed point.
\end{proof}

\section{The characteristic-zero DVR boundary}

The counterexample does not conflict with the special result cited in the AIM problem. Hacking proves the following in characteristic zero \cite[Proposition~10.14]{HackingThesis}.

\begin{proposition}[Hacking]
Let $T$ be the spectrum of a DVR and let $X\to T$ be a weakly $\mathbb Q$-Gorenstein family of semi-log-canonical surfaces. If the generic fiber is canonical, then the family is $\mathbb Q$-Gorenstein in the base-change sense.
\end{proposition}

A smoothing has smooth, hence canonical, generic fiber. Together with the general strong-to-weak implication, this yields equivalence in that restricted setting. The theorem does not cover arbitrary bases, nonreduced Artin thickenings, or positive characteristic.

There is also a bibliographic mismatch in the AIM report. Its bibliography names Hacking's later paper \cite{HackingPaper}, where the corresponding statement is Lemma~3.4, but the locator ``Proposition~10.14'' belongs to the earlier thesis \cite{HackingThesis}.

\section{Scope}

\begin{remark}[Scope]\label{rem:scope}
The base $B=\operatorname{Spec}k[t]/(t^2)$ is nonreduced, and the thickened total scheme is not normal. The fibers in the Lee--Nakayama family are of Fano or Del Pezzo type rather than canonically polarized. None of these properties invalidates the counterexample to the literal AIM question, which does not exclude them. Nevertheless, the example is not asserted to settle a modified problem restricted to stable or canonically polarized surfaces over reduced characteristic-zero bases.
\end{remark}

The proof above is a source-level reconstruction and requires no numerical computation.

\appendix

\section{Checklist for an independent referee}

\begin{enumerate}
  \item Read AIM Problem~1.1 and verify that it requires all $m>0$ and arbitrary base change without excluding nonreduced bases or positive characteristic.
  \item Check Lee--Nakayama Definition~7.1 and the weak/naive versus Koll\'ar terminology.
  \item Verify the fiber descriptions and flat projectivity in Example~7.4.
  \item Confirm that Remark~7.28 makes the first positive-characteristic Artin thickening naively $\mathbb Q$-Gorenstein.
  \item Confirm that Lemma~7.5 makes $\omega^{[2]}$ fail relative $S_2$ on the same thickening.
  \item Use Remark~7.6 to translate that failure into failure of the closed-fiber base-change map.
  \item Keep the unrestricted AIM question separate from Hacking's characteristic-zero DVR theorem and from canonically polarized-only variants.
\end{enumerate}

\begin{thebibliography}{9}

\bibitem{HackingThesis}
P.~Hacking,
\emph{A compactification of the space of plane curves},
Ph.D. thesis, University of Cambridge, arXiv:math/0104193v1 (2001).

\bibitem{HackingPaper}
P.~Hacking,
\emph{Compact moduli of plane curves},
Duke Math. J. \textbf{124} (2004), no.~2, 213--257;
arXiv:math/0310354v1.

\bibitem{LeeNakayama}
Y.~Lee and N.~Nakayama,
\emph{Grothendieck duality and $\mathbb Q$-Gorenstein morphisms},
Publ. Res. Inst. Math. Sci. \textbf{54} (2018), no.~3, 517--648.

\bibitem{VanOpstall}
M.~A. van Opstall,
\emph{Open problems in compact moduli spaces and birational geometry},
American Institute of Mathematics problem list, Problem~1.1.

\end{thebibliography}

\end{document}
